\section{Federated Query Containment Resolution}\label{sec:resolution}

In the previous section we defined the condition of query containment for federated queries however we did not touch on the algorithm to perform the containment evalution.
The goal of this section is to evaluate that.

The SPARQL specification~\cite{w3SPARQLFederated} define \texttt{SERVICE} the evaluation of the inner query without projection thus in the context of bag-set semantic the evaluation is the same as with set semantic~\cite{Afrati2010}.
Thus, for the evaluation the condition $B$ of ~\Cref{prop:condition-containment} any SPARQL query containment solver from the current literature
\footnote{See section \Cref{sec:related_works_qc} } can be orchestrated to implement condition $B$.
In that case the time complexity is $NP-complete$~\cite{Afrati2010}, like set-semantic query containmet.


With adjustement of the algorithm the complexity can be $O(n^2\log(n)\times)$~\cite{Afrati2010}, thus to satisfy condition $B$,
$O(n^2\log(n)\times n_{sQ1} \times n_{sQ2})$ where $n_{sQ1}$ and $ n_{sQ2}$ are respectively the number of service clauses from $Q_1$ and $Q_2$.
The adjustment consist of defining the bag-set query containment in terms of variables-onto containment mappings~\cite{Chaudhuri1993}.
Variable-onto mean that every variable in the sub query is present in the containment mapping.
% I need to define what i a containment mapping

The case is more complicated when considering \acl{BFR}.
Indeed without projection the problem cannot be reduced to set semantic query containment.
Thus, it is not possible to simply orchestrate existing set semantic query containment solvers.
However, it has the same time complexity then the bag-set semantic case~\cite{Afrati2010}.
To perform the computation of the containment it is necessary to consider goal-onto containment mappings~\cite{Chaudhuri1993, Afrati2010}.
Goal into containment mapping means that every cunjunct  as a containment mapping with the super query and that
the super query has more or the same number of cunjunct as the sub query.
% make it more formal and refer to the proposition from the paper (4.3)
